% !TeX program = pdflatex
\documentclass[12pt]{article}

% Idioma y codificación
\usepackage[spanish]{babel}
\usepackage[utf8]{inputenc}
\usepackage[T1]{fontenc}
\usepackage{lmodern}

% Página y espaciado (para acercarse a ~15 páginas)
\usepackage{geometry}
\geometry{letterpaper, margin=1in}
\usepackage{setspace}
\onehalfspacing

% Tipografía y microajustes
\usepackage{microtype}
\usepackage{csquotes}

% Tablas
\usepackage{booktabs}
\usepackage{longtable}
\usepackage{array}

% Listas
\usepackage{enumitem}

% Código
\usepackage{listingsutf8}
\usepackage{listings}
\lstset{inputencoding=utf8}
\lstset{literate=%
  {á}{{\'a}}1 {é}{{\'e}}1 {í}{{\'\i}}1 {ó}{{\'o}}1 {ú}{{\'u}}1
  {Á}{{\'A}}1 {É}{{\'E}}1 {Í}{{\'I}}1 {Ó}{{\'O}}1 {Ú}{{\'U}}1
  {ñ}{{\~n}}1 {Ñ}{{\~N}}1 {ü}{{\"u}}1 {Ü}{{\"U}}1
}
\usepackage{xcolor}
\lstdefinestyle{python}{
  language=Python,
  inputencoding=utf8,
  basicstyle=\ttfamily\small,
  numbers=left,
  numberstyle=\tiny,
  stepnumber=1,
  numbersep=8pt,
  breaklines=true,
  showstringspaces=false,
  tabsize=2,
  frame=single
}
\lstdefinestyle{sql}{
  language=SQL,
  inputencoding=utf8,
  basicstyle=\ttfamily\small,
  numbers=left,
  numberstyle=\tiny,
  stepnumber=1,
  numbersep=8pt,
  breaklines=true,
  showstringspaces=false,
  tabsize=2,
  frame=single
}

% Hipervínculos
\usepackage{hyperref}
\hypersetup{
  colorlinks=true,
  linkcolor=blue,
  urlcolor=blue,
  citecolor=blue
}

% Formato académico y paginación por sección
\usepackage{titlesec}
\usepackage{etoolbox}
\titlespacing*{\section}{0pt}{3ex plus 1ex minus .2ex}{2ex}
\titlespacing*{\subsection}{0pt}{2.5ex plus 0.8ex minus .2ex}{1.2ex}
\titlespacing*{\subsubsection}{0pt}{2ex plus 0.6ex minus .2ex}{1ex}
\preto\section{\clearpage} % cada sección inicia en nueva hoja

% Encabezados/pies limpios
\usepackage{fancyhdr}
\pagestyle{fancy}
\fancyhf{}
\lhead{Documento de apoyo (Sesión 2, Unidad 2)}
\rhead{Procesos ETL y Modelado OLAP}
\cfoot{\thepage}
\renewcommand{\headrulewidth}{0.4pt}
\renewcommand{\footrulewidth}{0pt}

% Bibliografía (biblatex + biber)
\usepackage[
  backend=biber,
  style=ieee,
  sorting=nyt,
  maxbibnames=99
]{biblatex}
\addbibresource{referencias_semana2.bib}

% Metadatos del documento
\title{\textbf{Documento de apoyo (Sesión 2, Unidad 2)}\\
\large Procesos ETL y Modelado de Datos para Análisis OLAP\\
\normalsize Diseño y optimización de cubos multidimensionales}
\author{Javier Andrés Valdez González / Auxiliar \\\\
\small Seminario de Sistemas 2 — Sección P \\
\small Universidad de San Carlos de Guatemala}
\date{Fecha: 04 de febrero de 2026}

\begin{document}
\maketitle

\begin{abstract}
Un cubo OLAP es una estructura lógica que permite analizar indicadores desde múltiples perspectivas (por ejemplo: tiempo, producto, región). Para que el análisis sea confiable y rápido, el cubo se alimenta de un modelo dimensional (hechos y dimensiones) construido a partir de procesos ETL/ELT: extracción, limpieza, transformación, integración y carga. Este documento describe los elementos de un cubo multidimensional, el diseño de dimensiones y jerarquías (incluyendo dimensiones normalizadas tipo copo de nieve), el manejo de cambios mediante Slowly Changing Dimensions (SCD), y la correcta definición de medidas y reglas de agregación. Además, compara los modelos dimensionales (estrella, copo de nieve y constelación) y resume un proceso práctico de diseño, junto con técnicas básicas de optimización como granularidad, agregaciones y particionamiento.
\end{abstract}

\section{Panorama general}
Un cubo OLAP es la estructura lógica que permite analizar indicadores desde múltiples perspectivas (por ejemplo: tiempo, producto, región). Para que el análisis sea confiable y rápido, el cubo se alimenta de un modelo dimensional (hechos y dimensiones) construido a partir de procesos ETL/ELT: extracción, limpieza, transformación, integración y carga.

\subsection*{Objetivos de la sesión}
\begin{itemize}[leftmargin=*, itemsep=0.25em]
  \item Comprender qué es un cubo multidimensional y cómo se relaciona con un modelo dimensional.
  \item Identificar dimensiones, jerarquías y medidas para responder preguntas de negocio.
  \item Comparar modelos estrella, copo de nieve y constelación para distintos escenarios.
  \item Aplicar principios básicos de optimización: granularidad, agregaciones y particiones.
\end{itemize}

\noindent\textbf{Fuentes:} \cite{kimball_toolkit, mslearn_cubes_multidimensional, mslearn_ssas_intro}.

\section{Definición de cubo multidimensional}
Un cubo multidimensional modela un conjunto de hechos (medidas numéricas) organizados por dimensiones. En la práctica, el ``cubo'' puede materializarse como un modelo tabular (columnar) o como un modelo multidimensional clásico; lo importante es que permita consultas por niveles (jerarquías) y agregaciones consistentes.

\subsection*{Elementos básicos}
\begin{itemize}[leftmargin=*, itemsep=0.25em]
  \item \textbf{Grupo de medidas:} conjunto de métricas relacionadas (p.\ ej., Ventas, Costos, Margen).
  \item \textbf{Dimensiones:} ejes de análisis (Tiempo, Producto, Cliente, Región, Canal).
  \item \textbf{Jerarquías:} niveles que permiten navegar de lo general a lo específico (Año \(\rightarrow\) Mes \(\rightarrow\) Día).
  \item \textbf{Granularidad:} nivel de detalle al que se registra el hecho (por ticket, por día, por línea de factura).
\end{itemize}

\noindent\textbf{Regla práctica:} antes de pensar en herramientas, define la pregunta de negocio y la granularidad. Esto evita modelos sobredimensionados o medidas incoherentes.

\noindent\textbf{Fuentes:} \cite{mslearn_cubes_multidimensional, mslearn_ssas_intro, kimballgroup_star_olap}.

\section{Dimensiones y jerarquías}
Las dimensiones describen el contexto de un hecho. Su diseño impacta directamente la experiencia de análisis: si la dimensión no refleja el lenguaje del negocio, el cubo será difícil de usar.

\subsection*{Buenas prácticas de diseño}
\begin{itemize}[leftmargin=*, itemsep=0.25em]
  \item Definir claves sustitutas (\textit{surrogate keys}) para dimensiones, separando la llave técnica de la llave de negocio.
  \item Mantener dimensiones conformadas (\textit{conformed}) cuando se reutilizan entre áreas (ventas, inventario, devoluciones).
  \item Modelar jerarquías naturales: Tiempo (Año \(\rightarrow\) Trimestre \(\rightarrow\) Mes \(\rightarrow\) Día), Geografía (País \(\rightarrow\) Depto.\ \(\rightarrow\) Ciudad).
  \item Evitar muchos atributos ``libres'' sin propósito: cada atributo debe soportar filtros o segmentaciones reales.
\end{itemize}

\subsection*{Jerarquías: ventajas}
\begin{itemize}[leftmargin=*, itemsep=0.25em]
  \item Permiten \textit{drill-down} y \textit{roll-up} sin redefinir reportes.
  \item Reducen errores al estandarizar niveles de agregación.
  \item Mejoran rendimiento cuando el motor usa agregaciones por nivel.
\end{itemize}

\noindent\textbf{Fuentes:} \cite{kimballgroup_star_olap, ibm_what_is_olap}.

\subsection*{Dimensiones \textit{snowflaked} (Normalizadas / Copo de Nieve)}
Una dimensión \textit{snowflaked} se normaliza en varias tablas relacionadas (subdimensiones). Es típica del esquema copo de nieve.

\subsubsection*{Ejemplo}
Dimensión Producto separada en: \textit{Producto}, \textit{Subcategoría}, \textit{Categoría} y \textit{Marca}.

\subsubsection*{Ventajas}
\begin{itemize}[leftmargin=*, itemsep=0.25em]
  \item Menor redundancia de datos.
  \item Puede facilitar mantenimiento de catálogos muy jerárquicos.
\end{itemize}

\subsubsection*{Desventajas}
\begin{itemize}[leftmargin=*, itemsep=0.25em]
  \item Consultas más complejas (más \textit{joins}).
  \item A veces reduce desempeño frente al esquema estrella.
\end{itemize}

\section{Dimensiones de cambio lento (Slowly Changing Dimensions, SCD)}
Las SCD manejan cambios en atributos de dimensión a través del tiempo (por ejemplo, un cliente cambia de ciudad o un producto cambia de categoría). La elección del tipo SCD determina cómo se conserva el historial.

\subsection*{SCD Tipo 0 (Sin cambios)}
Los atributos no se actualizan. Se conserva el valor original siempre. \textbf{Uso:} valores ``congelados'' o definidos al alta.

\subsection*{SCD Tipo 1 (Sobrescritura)}
Se actualiza el atributo y se pierde el valor anterior (sin historial). \textbf{Uso:} correcciones de errores, datos que no requieren historial (ej.\ corregir ortografía).

\subsection*{SCD Tipo 2 (Histórico completo)}
Se inserta un nuevo registro de la dimensión cuando hay cambios y se conserva el anterior, usualmente con:
\begin{itemize}[leftmargin=*, itemsep=0.25em]
  \item FechaInicio / FechaFin
  \item IndicadorActual (\textit{CurrentFlag})
  \item Nueva \textit{surrogate key}
\end{itemize}
\noindent\textbf{Uso:} cuando se necesita analizar ``como era'' el cliente/producto en el momento del evento.

\subsection*{SCD Tipo 3 (Historial limitado)}
Se guarda el valor anterior en otra columna (ej.\ \textit{CiudadActual} y \textit{CiudadAnterior}). \textbf{Uso:} se requiere sólo ``último cambio'' o historial corto.

\section{Medidas y agregaciones}
Las medidas son valores numéricos que se resumen (agregan) al consultar el cubo. Elegir mal la función de agregación genera resultados incorrectos (por ejemplo, sumar porcentajes).

\subsection*{Tipos de medidas (según su comportamiento de agregación)}
\begin{table}[h]
\centering
\begin{tabular}{@{}p{3cm}p{5.5cm}p{5.5cm}@{}}
\toprule
\textbf{Tipo} & \textbf{Ejemplo} & \textbf{Agregación típica} \\
\midrule
Aditiva & Ventas, unidades & SUM \\
Semi-aditiva & Saldo de cuenta & SUM por algunas dimensiones; LAST por tiempo \\
No aditiva & Tasa, \% margen & Promedio ponderado / medida calculada \\
\bottomrule
\end{tabular}
\caption{Tipos de medidas y reglas típicas de agregación.}
\end{table}

\subsection*{Recomendaciones rápidas}
\begin{itemize}[leftmargin=*, itemsep=0.25em]
  \item Definir medidas base (atómicas) y derivar KPIs como medidas calculadas.
  \item Evitar duplicidad de medidas; documentar definición y unidad (Q, \$, \%, etc.).
\end{itemize}

\noindent\textbf{Fuentes:} \cite{kimballgroup_star_olap, mslearn_star_schema_powerbi}.

\section{Modelos dimensionales: estrella, copo de nieve y constelación}
El modelo dimensional organiza tablas para análisis. La elección depende del equilibrio entre simplicidad (uso) y normalización (mantenimiento).

\subsection*{Descripción breve}
\begin{itemize}[leftmargin=*, itemsep=0.25em]
  \item \textbf{Modelo estrella:} una tabla de hechos central conectada a dimensiones desnormalizadas. Suele ser el recomendado para autoservicio en BI (p.\ ej., Power BI) por su simplicidad y desempeño.
  \item \textbf{Copo de nieve:} dimensiones normalizadas en subtablas (más \textit{joins}), con menor redundancia pero mayor complejidad.
  \item \textbf{Constelación:} múltiples tablas de hechos que comparten dimensiones conformadas (útil cuando existen varios procesos: ventas, devoluciones, compras, etc.).
\end{itemize}

\subsection*{Comparativa rápida}
\begin{table}[h]
\centering
\begin{tabular}{@{}p{2.3cm}p{4.1cm}p{4.1cm}p{4.1cm}@{}}
\toprule
\textbf{Modelo} & \textbf{Ventaja} & \textbf{Desventaja} & \textbf{Cuándo usar} \\
\midrule
Estrella & Simple y rápido para consultar & Puede repetir atributos & Autoservicio, reporting, tabular \\
Copo de nieve & Menos redundancia & Más joins; más complejo & Catálogos grandes y jerarquías formales \\
Constelación & Reutiliza dimensiones entre varios hechos & Mayor diseño/validación & Múltiples procesos (ventas, devol.) \\
\bottomrule
\end{tabular}
\caption{Comparación entre estrella, copo de nieve y constelación.}
\end{table}

\noindent\textbf{Fuentes:} \cite{kimballgroup_star_olap, ibm_what_is_olap, kimball_bus_architecture}.

\section{Proceso recomendado para diseñar un cubo}
Una secuencia práctica (de negocio a modelo):
\begin{enumerate}[leftmargin=*, itemsep=0.25em]
  \item Definir preguntas de negocio y KPI (qué decisión se quiere apoyar).
  \item Elegir el proceso de negocio (ventas, inventario, devoluciones) y su granularidad.
  \item Diseñar la tabla de hechos (medidas atómicas) y asegurar claves a dimensiones.
  \item Diseñar dimensiones con atributos y jerarquías (conformadas cuando aplique).
  \item Probar con consultas y ajustar el modelo.
\end{enumerate}

\noindent\textbf{Fuentes:} \cite{kimballgroup_star_olap, mslearn_star_schema_powerbi}.

\section{Optimización de cubos multidimensionales}
La optimización busca reducir tiempos de refresco y mejorar la respuesta a consultas. En general, se gana más optimizando el modelo y los datos que ``forzando'' la herramienta.

\subsection*{Técnicas comunes}
\begin{itemize}[leftmargin=*, itemsep=0.25em]
  \item \textbf{Granularidad adecuada:} evita cargar detalle innecesario si el negocio consulta a nivel agregado.
  \item \textbf{Particionamiento:} separar hechos por periodos (p.\ ej., por año o mes) para refrescos incrementales.
  \item \textbf{Índices y claves:} claves enteras y ordenadas mejoran \textit{joins} y compresión.
  \item \textbf{Columnas y cardinalidad:} evitar columnas de alta cardinalidad como texto libre en dimensiones.
  \item \textbf{Control de refresco:} incremental cuando sea posible; validar cambios de dimensiones (SCD).
\end{itemize}

\subsection*{Señales de alerta (para revisar antes de ``subir hardware'')}
\begin{itemize}[leftmargin=*, itemsep=0.25em]
  \item Tablas de hechos con campos de texto o llaves no numéricas.
  \item Dimensiones con atributos duplicados, o sin jerarquías claras.
  \item Medidas no aditivas calculadas como SUM (resultados incorrectos).
\end{itemize}

\noindent\textbf{Fuentes:} \cite{kimballgroup_star_olap, mslearn_star_schema_powerbi, mslearn_tabular_performance}.

\section{Cierre y bibliografía}
Un buen cubo OLAP no se mide por la cantidad de tablas o columnas, sino por su capacidad de responder preguntas del negocio de forma consistente y con buen rendimiento. La base es un modelo dimensional claro (hechos + dimensiones), alimentado por ETL confiable y con medidas correctamente definidas.

\subsection*{Checklist rápido para validar tu diseño}
\begin{itemize}[leftmargin=*, itemsep=0.25em]
  \item ¿La granularidad del hecho está claramente definida y documentada?
  \item ¿Las dimensiones reflejan el lenguaje del negocio y tienen jerarquías útiles?
  \item ¿Cada medida tiene una definición (fórmula), unidad y regla de agregación?
  \item ¿El modelo es estrella (o justifica por qué no) y evita relaciones ambiguas?
\end{itemize}

\nocite{*}
\printbibliography

\end{document}
